\section{Results}
\label{sec:results}

\subsection{Foundation models decode four-class motor imagery below every
supervised comparator}

Table~\ref{tab:heldout} gives held-out performance on BCI~IV-2a. The ordering
is total: every supervised arm, including the weakest one we ran, outperforms
every foundation-model configuration. The best foundation model, CBraMod
fine-tuned with its validation-selected recipe, reaches 0.4796 against a
0.5305 floor for the supervised arms and 0.7145 for the strongest. Against
each supervised comparator the paired accuracy comparison sits at the
attainable-$p$ floor with no subject wins.

\begin{table}[t]
\centering
\caption{Held-out performance, BCI~IV-2a, $n=9$, chance accuracy 0.250.
Condensed from the full grid of 26 configurations.}
\label{tab:heldout}
\small
\begin{tabular}{llccc}
\toprule
Configuration & Role & Acc & Brier & ECE \\
\midrule
seed ensemble            & supervised & 0.7145 & 0.3852 & 0.0829 \\
ShallowConvNet           & supervised & 0.6962 & 0.4047 & 0.0808 \\
ATCNet, broadband        & supervised & 0.6562 & 0.4695 & 0.1062 \\
EEG Conformer, broadband & supervised & 0.5305 & 0.5858 & 0.1072 \\
\midrule
CBraMod, fine-tuned      & FM         & 0.4796 & 0.6493 & 0.1236 \\
CBraMod, frozen          & FM         & 0.3927 & 0.7080 & 0.1030 \\
CBraMod, random init     & control    & 0.3407 & 0.7396 & 0.0993 \\
LaBraM, frozen           & FM         & 0.3322 & 0.7323 & 0.0874 \\
\bottomrule
\end{tabular}
\end{table}

\subsection{It is not a fine-tuning recipe artifact}
Under the reference recipe the best selection-Brier epoch was 0--4 of 60 across
all runs while warm-up alone lasts six, so the model peaked before the learning
rate reached its target --- what overfitting a 5.8\,M-parameter transformer to
230 trials looks like. A six-configuration sweep ranked on validation improves
fine-tuning substantially and leaves the ordering in Table~\ref{tab:heldout}
intact. Both remedies that work, freezing lower blocks and lowering the learning
rate, reduce how far pretrained weights can move, consistent with that
reading.

\subsection{The matched-input control changes sign across architectures}
\label{sec:decomp}

Because the two arms are fed different preprocessing by design
(Section~\ref{sec:twopaths}), a reader is entitled to ask whether the gap is
the models or the pipeline. The standard answer is a matched-input control:
retrain the supervised comparator on the arrays the foundation models actually
receive. Table~\ref{tab:decomp} runs that control on three architectures.

\begin{table}[t]
\centering
\caption{Matched-input control, BCI~IV-2a, $n=9$. The pipeline term is
narrowband minus broadband accuracy for the same architecture, so a negative
value means the foundation models' preprocessing \emph{helped} that
architecture. Wins counts subjects for which the narrowband input was better.
$p_{\mathrm{BH}}$ is corrected across every component test this decomposition
computes.}
\label{tab:decomp}
\small
\begin{tabular}{lccccc}
\toprule
Architecture & Narrow & Broad & Pipeline term & Wins & $p_{\mathrm{BH}}$ \\
\midrule
ATCNet         & 0.5788 & 0.6562 & $-0.078$ (helps) & 1/9 & 0.073 \\
ShallowConvNet & 0.6453 & 0.6084 & $+0.037$ (costs) & 5/9 & 0.304 \\
EEG Conformer  & 0.6182 & 0.5305 & $+0.088$ (costs) & 8/9 & 0.176 \\
\bottomrule
\end{tabular}
\end{table}

For ATCNet the foundation models' own preprocessing made the supervised
baseline \emph{better} by 0.078 accuracy, on 8 of 9 subjects, so for that
comparator there is no preprocessing explanation for the gap at all. For EEG
Conformer the same pipeline cost 0.088, on 8 of 9 subjects in the opposite
direction. The point estimates therefore differ in sign, and the subject-level
directions are opposed.

We state this carefully. No individual pipeline term clears
Benjamini--Hochberg correction across this decomposition's family, so none is
separately established at $n=9$. What the table supports is narrower and
sufficient: the measured term is not stable across comparator architectures and
is not reliably estimable from any one of them, so a single-architecture
matched-input control is not a dependable basis for splitting a
pretrained-versus-supervised gap between model and pipeline --- which is how
such controls are ordinarily used.

The term is also not a filter-band term: the pipelines differ in band, filter
design, sampling rate and notching, and the broadband arm receives its own
per-subject architecture search, so the column measures the pipeline together
with the retuning it requires. We report raw gaps rather than a percentage
split, because a share is only interpretable when its denominator has a stable
sign. Nothing in the argument is specific to EEG, and to our knowledge the
contrast has not previously been reported.

\subsection{The deficit is task-dependent}
On two-class BNCI2014\_004 the picture changes. CBraMod fine-tuned reaches
0.8085 against 0.8458 for the best supervised arm, 0.8185 for a tuned
ShallowConvNet and 0.7981 for EEG Conformer --- competitive, and above one of
the supervised comparators. The four-class deficit in Table~\ref{tab:heldout}
is therefore not a general statement about these models on motor imagery: we
could not detect it here. We do not attribute the difference to class count
alone. BNCI2014\_004 differs from BCI~IV-2a in dataset, in class count, and in
carrying three electrodes rather than 22, and this design does not separate
those factors.

\subsection{Robustness}
Two checks, both requested by the obvious objections.

\emph{Seeds.} The decisive comparison --- CBraMod fine-tuned against its
architecture-matched random-init control --- was repeated at three independent
reseeds of the entire pipeline, which redraws the fit/selection split, the
initialisation and the batch ordering. On BCI~IV-2a the subject-averaged effect
is 0.139, 0.132 and 0.130; on BNCI2014\_004 it is 0.107, 0.111 and 0.108. Every
one of the six runs is individually significant, and the lowest subject win
rate anywhere is 8 of 9. We do not quote these as a mean with a standard
deviation: that spread measures reproducibility under training randomness, not
uncertainty about the effect, which is governed by $n=9$ subjects.

\emph{Budget.} A reader may object that the random-init arm was simply
undertrained. Repeating the BCI~IV-2a comparison at 240 epochs, four times the
budget, for both arms equally: random initialisation gained 0.017 accuracy and
closed 0.024 of the gap, while the pretrained arm gained nothing. At the larger
budget the effect is significant on all nine subjects. The random-init arm reached its best
validation Brier at epoch 218 of 240 with the curve flat from about 190, so it
had plateaued rather than been truncated --- and its best value, 0.755, is
essentially the four-class chance Brier of 0.750.

\subsection{Calibration is where they hold up}
The raw logits are badly overconfident: frozen configurations require
temperatures of 3.5 to 9.5, with one to four of nine subjects saturating the
clamp at 20.0, against 1.03 for a tuned supervised convnet. After that single
validation-fitted scalar, held-out ECE lands inside the supervised range ---
0.103 to 0.124 for CBraMod against 0.101 to 0.143 across the tuned supervised
arms --- despite far lower accuracy.
